problem:  
Let \( X \) be a smooth Fano manifold and \( \omega_t = \omega_0 + i\partial\bar{\partial}\varphi_t \) evolve by the normalized Kähler–Ricci flow  
\[
\frac{\partial \varphi_t}{\partial t} = \log \frac{(\omega_0 + i\partial\bar{\partial}\varphi_t)^n}{\omega_0^n} + \varphi_t - h_0, 
\qquad \int_X e^{-\varphi_t}\omega_0^n = V,
\]
for a fixed reference potential \( h_0 \).  
For each \(t > 0\) and integer \(m > 0\), consider the *multiplier ideal sheaf*  
\[
\mathcal{I}_t^{(m)} := \mathcal{J}(m\varphi_t),
\]
associated to the weight \( m \varphi_t \).  This defines a family of coherent ideals in \( \mathcal{O}_X \) that reflect the singularity type of the evolving potentials along the flow.  

For each \(m>0\), let
\[
I_m(t) := \mathcal{I}_t^{(m)},
\qquad
R = \bigoplus_{k\ge 0} H^0(X, -kK_X),
\]
and denote by \(\mathcal{F}^t R\) the *filtration of the anticanonical ring* generated by these ideals:
\[
\mathcal{F}^t R_k := H^0(X, -kK_X\otimes \mathcal{I}_t^{(k)}).
\]

Assume that \(\{\varphi_t\}\) remains uniformly bounded in \(C^0\) and that the family of filtrations \(\mathcal{F}^t R\) subconverges (in the Boucksom–Hisamoto–Jonsson sense) to a non‑Archimedean metric \(\phi_\infty^{\mathrm{NA}}\) on \(-K_X\).  

**Question:**  
Does \(\phi_\infty^{\mathrm{NA}}\) necessarily minimize the non‑Archimedean Ding functional \(\mathcal{D}^{\mathrm{NA}}\) over all non‑Archimedean metrics on \(-K_X\)?  
Equivalently, does the family of multiplier ideals \(\{\mathcal{I}_t^{(m)}\}\) extracted from the Kähler–Ricci flow encode the *optimal destabilizing test configuration* of \(X\)?  
In the K‑polystable case, must \(\phi_\infty^{\mathrm{NA}}\) coincide with the trivial non‑Archimedean metric arising from the smooth Kähler–Einstein metric?  

Why is it a "good" problem:  
This version uses mathematically robust objects—multiplier ideals and filtrations defined by sections of powers of the anticanonical bundle—whose asymptotic behavior is central in both complex analytic and algebraic formulations of K‑stability. It proposes a precise mechanism by which the analytic dynamics of the Kähler–Ricci flow could be translated into the algebraic language of filtrations and non‑Archimedean metrics.  

Resolving this problem would establish a canonical analytic construction of the valuations or filtrations predicted by stability theory and answer whether the Kähler–Ricci flow asymptotically detects the test configuration realizing the minimal non‑Archimedean Ding energy. The question lies at the heart of the analytic side of the Yau–Tian–Donaldson conjecture and the Boucksom–Jonsson theory of degeneration, offering a rigorous and deeply integrative bridge between nonlinear parabolic PDEs, pluripotential theory, and algebraic geometry.